\documentclass[12pt,a4paper]{article}
% --------------------
% Packages
% --------------------
\usepackage[a4paper,margin=1in]{geometry}
\usepackage{graphicx}
\usepackage{setspace}
\usepackage{lmodern}
\usepackage{amsmath}
\usepackage{float}
\usepackage{tikz}
\usetikzlibrary{shapes.geometric, arrows, positioning}
\usepackage{booktabs}
\usepackage{listings}
\usepackage{xcolor}
\usepackage{hyperref}

\lstset{
  basicstyle=\ttfamily\small,
  breaklines=true,
  frame=single,
  backgroundcolor=\color{gray!10},
  keywordstyle=\color{blue},
  commentstyle=\color{green!50!black},
}

\begin{document}
\thispagestyle{empty}
% --------------------
% LOGO (TOP CENTER)
% --------------------
\begin{center}
    
\includegraphics[width=15cm]{logo_deib.png}
\end{center}
\vspace{1.5cm}
% --------------------
% TITLE BLOCK
% --------------------
\begin{center}
    {\Huge\bfseries Internet Of Things\\[0.5cm]
    {\Large Challenge 2 -- Packet Sniffing}}
\end{center}
\vspace{2cm}
% --------------------
% AUTHOR INFORMATION
% --------------------
\begin{center}
    {\Large * Seyedaman Zargari -- 11109586} \\[0.3cm]
    {\Large Amirhossein Dibaj -- 11109540}
\end{center}
% --------------------
% PUSH CONTENT UP, YEAR DOWN
% --------------------
\vfill
% --------------------
% ACADEMIC YEAR (BOTTOM)
% --------------------
\begin{center}
    {\Large Academic Year 2025--2026}
\end{center}
\newpage
\tableofcontents
\thispagestyle{empty}
\newpage

\section{CQ1 -- NON DELETE Requests to coap.me}

\subsection{CQ1a}

\textbf{Filter:} All NON-Confirmable (type=1) DELETE (code=4) requests directed to
\texttt{coap.me} (IP \texttt{134.102.218.18}, resolved via DNS in the capture),
matched with their responses by token.

Out of 20 NON DELETE requests to \texttt{coap.me}, only \textbf{1} received a
successful (2.xx) response:

\begin{table}[h]
\centering
\begin{tabular}{@{}llll@{}}
\toprule
MID & Resource & Response Code & Meaning \\
\midrule
30800 & /validate & 2.02 & Deleted \\
\bottomrule
\end{tabular}
\caption{NON DELETE requests with successful response}
\end{table}

The remaining 19 requests all received \texttt{4.05 Method Not Allowed}.

\textbf{Answer CQ1a:} 1 request. \textbf{MID = 30800}.

\subsection{CQ1b}

Although the server responded with \texttt{2.02 Deleted} to MID 30800, a subsequent
GET request to \texttt{/validate} (frame 11363) still received a \texttt{2.05 Content}
response (frame 11364), proving the resource was not actually deleted.
\texttt{coap.me} is a public test server that acknowledges the DELETE method on
\texttt{/validate} but does not truly remove the resource.

\textbf{Answer CQ1b: 0}

\section{CQ2 -- Resources with Equal Unsuccessful CON POST and PUT}

\textbf{Method:} Extracted all Confirmable (type=0) POST (code=2) and PUT (code=3) requests
directed to the local server (\texttt{127.0.0.1:5683}). Matched each with its response
(by token, from source port 5683). Counted requests as unsuccessful only when a
4.xx or 5.xx response was present. Used MID as uniqueness criterion (retransmissions
share the same MID).

\begin{table}[h]
\centering
\begin{tabular}{@{}lccc@{}}
\toprule
Resource & Unique Unsucc. POST (X) & Unique Unsucc. PUT (Y) & X=Y? \\
\midrule
/basic                      & 2 & 1 & No  \\
/dining\_room               & 1 & 1 & \textbf{Yes} \\
/dining\_room/door          & 2 & 3 & No  \\
/dining\_room/temperature   & 5 & 2 & No  \\
/hello\_post                & 3 & 0 & No  \\
/hello\_world               & 4 & 0 & No  \\
/living\_room               & 1 & 0 & No  \\
/living\_room/door          & 0 & 1 & No  \\
/living\_room/temperature   & 1 & 0 & No  \\
/main\_door                 & 1 & 2 & No  \\
/test                       & 2 & 0 & No  \\
\bottomrule
\end{tabular}
\caption{Unique unsuccessful CON POST/PUT per resource on local server}
\end{table}

Only \texttt{/dining\_room} has $X = Y = 1$ (POST got 4.00 Bad Request; PUT got 4.04 Not Found).

\textbf{Answer CQ2: 1}

\section{CQ3 -- Observe Notifications for /dining\_room/temperature}

\subsection{CQ3a}

\textbf{Method:} Identified the observe registration (frame 4343: CON GET with
\texttt{Observe=0}, token \texttt{89b9}, to \texttt{127.0.0.1:5683}). Then collected all
responses from the server (srcport 5683) matching this token with the Observe option
present. Grouped by MID to eliminate retransmissions.

\begin{table}[H]
\centering
\begin{tabular}{@{}cllcccc@{}}
\toprule
\# & MID & Type & Observe & Temp (\textdegree C) & ACKed & Retransmissions \\
\midrule
1  & 10588 & ACK & 1  & 19.97 & Yes & 0 \\
2  & 35    & CON & 1  & 22.32 & Yes & 0 \\
3  & 38    & CON & 2  & 18.85 & Yes & 0 \\
4  & 43    & CON & 3  & 25.98 & Yes & 0 \\
5  & 45    & CON & 4  & 15.59 & Yes & 0 \\
6  & 49    & CON & 5  & 25.35 & Yes & 0 \\
7  & 54    & CON & 6  & 26.21 & \textbf{No} & 7 \\
8  & 62    & CON & 7  & 27.53 & \textbf{No} & 3 \\
9  & 74    & CON & 8  & 26.98 & \textbf{No} & 1 \\
10 & 87    & CON & 9  & 25.32 & \textbf{No} & 1 \\
11 & 90    & CON & 10 & 19.57 & \textbf{No} & 1 \\
\bottomrule
\end{tabular}
\caption{Unique observe notifications for \texttt{/dining\_room/temperature}}
\end{table}

\textbf{Answer CQ3a: 11}

\subsection{CQ3b}

Starting from notification \#7 (MID 54, Observe=6), the client stops acknowledging
CON notifications. MID~54 is retransmitted 7 times without receiving an ACK. According
to RFC~7641 \S4.5, when a CON notification fails (exhausts retransmissions without ACK),
the server should cancel the observe relationship. Despite this, the server continues
sending notifications \#8--11 (MIDs 62, 74, 87, 90), none of which are ever
acknowledged. All 5 unacknowledged notifications (MIDs 54, 62, 74, 87, 90) represent
wasted traffic.

\textbf{Answer CQ3b: 5}

\section{CQ4 -- MQTT-SN Messages from Local Broker (Port 1885)}

\textbf{Filter:} Configured tshark to decode UDP port 1885 as MQTT-SN
(\texttt{-d udp.port==1885,mqttsn}). Filtered for packets with
\texttt{udp.srcport == 1885} (messages FROM the broker to clients).

\textbf{Analysis:} 219 MQTT-SN messages were sent by clients \emph{to} port 1885
(\texttt{udp.dstport == 1885}). Every single one triggered an ICMP
``Destination Unreachable -- Port Unreachable'' (Type~3, Code~3) response, confirming
that no broker was listening on port 1885. Zero packets were observed with
\texttt{udp.srcport == 1885}, meaning the broker never sent any response.

\textbf{Answer CQ4: 0}

\section{CQ5 -- Last Will via Wildcard Subscription}

\textbf{Method:} Identified 4 CONNECT packets with Last Will flags set. Verified that all 4 clients
disconnected ungracefully (no DISCONNECT message found in their TCP streams), meaning all 4 last wills
were triggered. Then matched each will topic against all wildcard SUBSCRIBE topics on the same broker.

\begin{table}[H]
\centering
\begin{tabular}{@{}llll@{}}
\toprule
Will Topic & Broker & Matching Wildcard Sub & Subscriber Stream \\
\midrule
university/department12/room1/temperature & local (::1) & \texttt{university/\#} & 10 \\
metaverse/room2/floor4 & mosquitto & (none) & -- \\
hospital/facility3/area3 & mosquitto & (none) & -- \\
metaverse/room2/room2 & mosquitto & (none) & -- \\
\bottomrule
\end{tabular}
\caption{Last will topics and wildcard subscription matches}
\end{table}

Only the local broker will topic is matched by a wildcard subscription (\texttt{university/\#} on stream~10).

\textbf{Answer CQ5: 1}

\section{CQ6 -- Erasing Retained Values on HiveMQ}

\subsection{CQ6a}

\textbf{Method:} Filtered for PUBLISH messages to HiveMQ (IPs \texttt{18.192.151.104},
\texttt{35.158.43.69}) with \texttt{retain=1} and empty payload (\texttt{payload\_length = 0}).
An empty retained publish erases the previously stored value for that topic.

Three distinct TCP streams (clients) sent such messages: streams 1, 22, and 43.

\textbf{Answer CQ6a: 3}

\subsection{CQ6b}

\textbf{Method:} Extracted the Client ID from the CONNECT packet of each stream identified in CQ6a.

\begin{table}[H]
\centering
\begin{tabular}{@{}llcl@{}}
\toprule
Stream & Client ID & Length (bytes) & $>7$? \\
\midrule
1  & pcbtjyj    & 7  & No \\
22 & giuxfijwus & 10 & \textbf{Yes} \\
43 & atqeblw    & 7  & No \\
\bottomrule
\end{tabular}
\caption{Client IDs of clients erasing retained values}
\end{table}

\textbf{Answer CQ6b: 1}

\section{CQ7 -- Subscribe Requests with Multiple Wildcards}

\textbf{Filter:} SUBSCRIBE packets (\texttt{mqtt.msgtype == 8}) directed to the local broker
(\texttt{ipv6.dst == ::1}), where at least one topic filter contains two or more
wildcard characters (\texttt{+} or \texttt{\#}).

Nine such SUBSCRIBE requests were found:

\begin{table}[H]
\centering
\begin{tabular}{@{}cl@{}}
\toprule
Frame & Topic Filter \\
\midrule
226  & university/+/+/deposit \\
363  & hospital/+/+/light \\
460  & metaverse/+/room1/\# \\
473  & house/department3/+/+ \\
700  & metaverse/facility0/+/+ \\
1177 & hospital/fxlugs/+/\# \\
1915 & metaverse/+/+/hydraulic\_valve \\
2360 & factory/+/\# \\
2994 & house/+/room1/\# \\
\bottomrule
\end{tabular}
\caption{Subscribe requests with $\geq 2$ wildcards to local broker}
\end{table}

\textbf{Answer CQ7: 9}

\section{CQ8 -- Topic Depth Histogram (Both Files)}

\textbf{Method:} Extracted all MQTT PUBLISH messages (\texttt{mqtt.msgtype == 3}) directed to the
local broker in each capture:
\begin{itemize}
    \item \textbf{A.pcapng:} local broker at \texttt{127.0.0.1:1883}
    \item \textbf{B.pcapng:} local broker at \texttt{[::1]:1883}
\end{itemize}

For each message, the number of topic layers was computed by counting segments separated by
\texttt{/} (e.g., \texttt{university/room2/floor4} $\rightarrow$ 3 layers).

\begin{table}[H]
\centering
\begin{tabular}{@{}ccc@{}}
\toprule
Number of Layers & A.pcapng & B.pcapng \\
\midrule
1 & 0   & 21  \\
2 & 154 & 172 \\
3 & 131 & 162 \\
4 & 154 & 179 \\
\bottomrule
\end{tabular}
\caption{Topic layer distribution for PUBLISH messages to local broker}
\end{table}

\textbf{CQ8a} (total publishes in A.pcapng): \textbf{439}

\textbf{CQ8b} (total publishes in B.pcapng): \textbf{534}

\begin{figure}[H]
    \centering
    \includegraphics[width=0.85\textwidth]{cq8_histogram.png}
    \caption{Distribution of MQTT topic depth for local broker}
    \label{fig:cq8}
\end{figure}

B.pcapng shows a broader distribution with topics at all depths 1--4, including 21 single-layer topics
not present in A.pcapng. Both captures concentrate most traffic at 2--4 layers, with B.pcapng
having a consistently higher message count across all depths.

\end{document}